%-----------------------------------------------------------------
% BAB II TINJAUAN KEPUSTAKAAN
%
% Struktur sesuai panduan penulisan tesis:
%   2.1 Teori A
%   2.2 Teori B
%   2.3 Teori C
%   2.4 Teori D
%
% Jumlah subbab disesuaikan dengan kebutuhan; ganti nama subbab
% dengan topik kajian Anda dan tambahkan bila diperlukan.
%-----------------------------------------------------------------
\chapter{TINJAUAN KEPUSTAKAAN}

% Bab ini cukup memuat judul besar tiap topik kajian beserta
% uraiannya. Rujukan diambil dari `daftar-pustaka.bib`.

\section{[Teori A]}

\section{[Teori B]}

\section{[Teori C]}

\section{[Teori D]}

% \section{[Metode atau Konsep Pendukung]}

% \section{Penelitian Terkait}